\section{Threats to Validity}
\label{sec:threats}

\paragraph{Internal validity.}
The measurement reports per-method clause-count diffs. These are
deterministic given the inputs: rerunning the harness produces
byte-identical metric output (we verified). The 5{,}763 / 2{,}001
new-precondition counts are therefore exact for the inputs reported.
The shape categorisation uses a simple regex-based bucketing; we
documented its order-of-precedence (structural shape over content) so
the buckets are reproducible. A small number of clauses in the
\emph{other} bucket (1.5--1.7\%) are arithmetic comparisons our
regex does not match (e.g.\ \texttt{abs(x) <= K}); their existence
does not affect the headline result.

\paragraph{External validity.}
The measurement is on two libraries (Commons Lang and Commons IO),
not on a broader corpus. Generalisation to less utility-shaped code
(application logic, framework code, generics-heavy collections) is
left to follow-up. The exclusion of Guava is forced by a stack-depth
limit in the JavaParser-driven harness we use to build the
specification cache (the harness recurses on large compilation units;
raising \texttt{-Xss} would let Guava run, but our default
configuration is the realistic one for typical Maven test
invocations). The fix is mechanical (either raise the stack size or
rewrite the call-graph builder iteratively); we do not expect Guava
to qualitatively change the conclusion, because the structural classes
of additions (branch-conditional, polymorphic-dispatch) are
properties of the language, not the library.

The single-pass baseline against which we compare is the inferrer's
\texttt{InterproceduralAnalyzer} from~\cite{edmonds2026inference}.
A different single-pass implementation might propagate more
information per call site (e.g.\ guarded propagation in a single
pass), narrowing the gap. We do not claim the second pass is
strictly necessary against all conceivable single-pass strategies;
we claim it materially extends the specific baseline of an
established prior system.

\paragraph{Construct validity.}
``Refined'' is operationalised as ``the spec gained at least one
clause that was not in the baseline'', using exact string equality
on clause text. Two clauses that are semantically equivalent but
syntactically different (e.g.\ \texttt{a < b} vs.\ \texttt{b > a})
would count as additions even if redundant. The analyser does not
normalise clauses before equality, so the new-clause count is an
upper bound on the number of genuinely new \emph{logical}
obligations. The shape categorisation is unaffected: a duplicate
\texttt{a < b} is still a range/bound clause.

The callee-resolution step uses simple-name matching, which can
return a wrong callee for overloaded methods. The
\texttt{substituteArguments} step then rejects the substitution if
the formal and actual arity disagree; in our run, 4{,}631 of 7{,}719
Commons Lang call sites pass the arity check (60\%). The remaining
40\% are either inferrer-internal calls without a cached callee or
name collisions with the wrong callee. A full type-resolving
harness would close this measurement gap, but the affected call
sites would either add additional clauses (strengthening the
headline result) or be rejected by arity (no effect on the result).
The reported numbers are therefore a \emph{lower bound} on the
true compositional gain.

\paragraph{Conclusion validity.}
The empirical claims here are exact measurements, not statistical
inferences; no significance test is meaningful. The conclusions
(``the second pass refines 43.85\% of methods on Commons Lang'',
``branch-conditional implications dominate the additions'') hold for
the specific corpora reported and for the specific
\texttt{CompositionalAnalyzer} implementation in the artefact. We do
not claim that the same percentages would obtain on other libraries
or on a different second-pass implementation; the structural
explanation (single-pass propagation cannot lift through guards or
dispatch) suggests the qualitative finding will transfer, but
quantitative transfer is an open question.
